\section{Stability under displacement and splitting of zeros}
\label{sec:stability}

Fix the labelled character data, a matrix norm bound $M$, an erasure budget $K$ and an insertion budget $\ell$. Write $m$ for the top multi-index, and let $d$ be the number of proper multi-indices $\alpha\ne m$ in the expansion
\[
 P_A=E_m+\sum_{\alpha\ne m} c_\alpha(A)E_\alpha.
\]
Use ordinary Euclidean norms on vectors of polynomial coefficients in increasing monomial order. The space $V_\ell$ consists of $\sum_{\alpha\ne m} h_\alpha E_\alpha$ with $\deg h_\alpha\le\ell$. Its coefficient norm is well-defined by rational linear independence. Let $T_0$ be a uniform window on which every nonzero member of $V_\ell$ misses more than $K$ zero multiplicities of every $P_A$ in the matrix ball, and put $T=T_0+1$. The margin permits roots at $\pm T$ to enter or leave the observation window as $A$ varies.

For a finite real multiset $D$ with $b$ elements counted with multiplicity, put $p_D(z)=\prod_{x\in D}(z-x)$. Write $I_D f$ for the unique polynomial of degree less than $b$ whose difference from the entire function $f$ is divisible by $p_D$ in the ring of entire functions. Thus a point of multiplicity $v$ prescribes derivatives of orders $0,\ldots,v-1$. Set $I_\varnothing f=0$. The coefficient norm of $I_Df$ treats nearby or coinciding nodes continuously.

\begin{lemma}\label{lem:stableinterpolation}
Fix a bound $B$ on $b$, and restrict all nodes to $[-U,U]$. If $D,D'$ have the same size and can be paired with distance at most $\delta$, then
\[
 \|I_Df-I_{D'}f\|\le C\delta\max_{|w|=U+2}|f(w)| .
\]
If $D\subseteq Y$ as multisets with $|Y|\le B$, then $\|I_D f\|\le C'\|I_Y f\|$. Both constants depend only on $B,U$.
\end{lemma}

\begin{proof}
On the positively oriented circle $|w|=U+2$, Cauchy's formula gives
\[
 I_D f(z)=\frac{1}{2\pi i}\int_{|w|=U+2}
 f(w)\frac{p_D(w)-p_D(z)}{(w-z)p_D(w)}\,dw .
\]
The right side is a polynomial in $z$ of degree less than $b$, and subtraction from $f(z)$ leaves a holomorphic multiple of $p_D(z)$ near the nodes. Each coefficient of its kernel has bounded first derivatives in the $b$ node variables on the compact box, since $|p_D(w)|\ge 2^b$. Taking maxima for $b\le B$ proves the first assertion, including collisions. For the second, $I_Df=I_D(I_Yf)$ follows from divisibility. The displayed integral is bounded uniformly on polynomials of degree less than $B$ with unit coefficient norm.
\end{proof}

\begin{lemma}\label{lem:stablefloor}
There is $a>0$ such that for all $\|A\|\le M$, every multiset $Z$ obtained from the zeros of $P_A$ in $[-T,T]$ by at most $K$ deletions, and every $H=\sum_{\alpha\ne m}h_\alpha E_\alpha\in V_\ell$, one has
\[
 \|I_ZH\|\ge a\,\|(h_\alpha)_{\alpha\ne m}\|.
\]
\end{lemma}

\begin{proof}
The number of zeros of $P_A$ in $[-T,T]$ is uniformly bounded. Indeed, local uniform convergence along any convergent subsequence of matrices gives a nonzero limit $P_{A_*}$. A circle enclosing the interval and avoiding its zeros then gives a common bound by Rouch\'e's theorem. If the claimed inequality failed, take a sequence $A_j\to A_*$, unit coefficient vectors converging to a unit vector defining $H_*$, and multisets $Z_j$ with $I_{Z_j}H_j\to0$. Pass to a subsequence of common cardinality with all ordered nodes convergent to a multiset $Z_*$. The integral formula gives $I_{Z_*}H_*=0$. Around each zero of $P_{A_*}$ in $[-T_0,T_0]$ take a disjoint small disk with its real diameter inside $(-T,T)$ and with boundary avoiding zeros of $P_{A_*}$ and points of $Z_*$. The exact zero counts of $P_{A_j}$ in the disks equal those of $P_{A_*}$ for large $j$, and all these zeros are real. Shrink the disks to exclude all points of $Z_*$ other than their respective centers. At most $K$ labels across these disks were deleted, so $Z_*$ contains all but at most $K$ multiplicities of the inner-window zeros of $P_{A_*}$. The divisibility $p_{Z_*}\mid H_*$ contradicts the choice of $T_0$.
\end{proof}

\begin{theorem}\label{thm:stabledecoder}
Suppose $\|A\|\le M$ and $0\le\delta\le1$. Erase at most $K$ labels of the zeros of $P_A$ in $[-T,T]$, displace each retained label by at most $\delta$, and insert at most $\ell$ arbitrary labels in $[-T-1,T+1]$. Let $Y$ be the resulting multiset, with no information about the matching. There are constants $R_0,C$ depending only on the fixed character list, $M,K,\ell,T$, such that the following estimator satisfies
\[
 \big\|(\widehat c_\alpha-c_\alpha(A))_{\alpha\ne m}\big\|\le C\delta .
\]
Choose any minimizer of
\[
 \left\| I_Y\!\left(W E_m+\sum_{\alpha\ne m}Q_\alpha E_\alpha\right)\right\|
\]
over monic $W$ of degree $\ell$ and $\deg Q_\alpha\le\ell$, with the joint coefficient norm bounded by $R_0$. Return $\widehat c_\alpha=[z^\ell]Q_\alpha$. The minimization is convex, with $\ell+d(\ell+1)$ free scalar coefficients. Repeated genuine zeros may split into separate reported locations, and distinct zeros may be reported at the same location.
\end{theorem}

\begin{proof}
Choose $R_0$ to contain all coefficient vectors
$(W_*,(c_\alpha(A)W_*)_\alpha)$ for $\|A\|\le M$ and all monic $W_*$ of degree $\ell$ whose roots lie in $[-T-1,T+1]$. Such a radius is finite by continuity and compactness. In particular, for the actual insertion multiset $J$, one may take $W_*=z^{\ell-|J|}p_J$. If $Z$ denotes the retained genuine labels before displacement, then $p_{Z\cup J}$ divides $W_*P_A$, where multiset union here adds multiplicities. The multiset $Z\cup J$ pairs with $Y$ at distance at most $\delta$. The first interpolation bound applied to the uniformly bounded family $W_*P_A$ shows that this feasible vector has residual at most $C_1\delta$. A minimizer exists by compactness and has at most this residual.

Write $G=W E_m+\sum Q_\alpha E_\alpha$ for its fitted function and $Z'$ for the displaced genuine submultiset of $Y$. The second interpolation bound gives $\|I_{Z'}G\|\le C_2\delta$. All fitted coefficient vectors lie in the same ball, so $|G|$ on the integration circle is uniformly bounded. The first bound then gives $\|I_ZG\|\le C_3\delta$. All cardinalities used here are bounded by the uniform count in the proof of the preceding lemma plus $\ell$. On $Z$ we have $I_Z(WP_A)=0$. Thus $H=G-WP_A=\sum (Q_\alpha-c_\alpha(A)W)E_\alpha$ satisfies the same residual bound there. Lemma~\ref{lem:stablefloor} yields
\[
 \big\|(Q_\alpha-c_\alpha(A)W)_{\alpha\ne m}\big\|\le (C_3/a)\delta .
\]
Since $W$ is monic, taking the coefficient of $z^\ell$ proves the result.
\end{proof}

The constants are existential. In particular, the theorem supplies no numerical condition number for a given list of characters. The contour representation makes the bound uniform over collisions; solving a system in unscaled point-evaluation rows can have very different numerical conditioning.
